\subsection{\texorpdfstring{\(K_{0}(A)\) 的换环}{K\_\{0\}(A) 的换环}}

设 A 与 B 是环。设 \(f: A \to B\) 是一个环同态。若 P 是一个有限生成投射 A-模，则 B-模 \(f^{*}(P) = B \otimes_{A} P\) 是投射且有限生成的（II, §5, No.~2, 第 281 页，推论）。映射 \(P \mapsto \text{cl}(f^{*}(P))\) 是从幺半群 \(\mathcal{P}(A)\) 到幺半群 \(\mathcal{P}(B)\) 的一个同态，因此定义了一个同态 \(f^{*}: K_{0}(A) \to K_{0}(B)\)，它由关系式 \(f^{*}([P]_{\mathcal{P}(A)}) = [f^{*}(P)]_{\mathcal{P}(B)}\)（对每个有限生成投射 A-模 P）刻画。若 \(g: B \to C\) 是第二个环同态，则由标量扩张的递移性（II, §5, No.~1, 第 278 页，命题 2）得出，同态 \((g \circ f)^*\) 与 \(g^{*} \circ f^{*}\)（都是从 \(\mathrm{K}_{0}(\mathrm{A})\) 到 \(\mathrm{K}_{0}(\mathrm{C})\) 的同态）相等。

假设 \(f\) 使 B 成为一个有限生成投射左 A-模。设 Q 是一个有限生成投射左 B-模。则 Q 是一个有限生成自由 B-模的直因子，而后者作为 A-模是投射且有限生成的。于是，由 Q 经标量限制得到的 A-模 \(f_{*}(\mathrm{Q})\) 是投射且有限生成的。同上，我们导出一个同态 \(f_{*}: \mathrm{K}_{0}(\mathrm{B}) \to \mathrm{K}_{0}(\mathrm{A})\)，它由关系式 \(f_{*}([Q]_{\mathcal{P}(B)}) = [f_{*}(Q)]_{\mathcal{P}(A)}\)（对每个有限生成投射 B-模 Q）刻画。若 \(g: \mathrm{B} \to \mathrm{C}\) 是使 C 成为有限生成投射 B-模的环同态，则同态 \((g \circ f)_{*}\) 与 \(f_{*} \circ g_{*}\)（都是从 \(\mathrm{K}_{0}(\mathrm{C})\) 到 \(\mathrm{K}_{0}(\mathrm{A})\) 的同态）相等。

